\section{Title, abstract, and main text}
Write the title in the file \texttt{title.tex}, the abstract in the file \texttt{abstract.tex}, and all other text in \texttt{body.tex}.
It is recommended to use one line per sentence in the latex files because it makes diffs between versions more readable.

\section{Draft versions}
Fill the information about the paper draft in the file \texttt{draft.tex} and about internal/conference notes in \texttt{note.tex}.

\section{Figures}
Put your figures in the subdirectory \texttt{figures}.
Add them to the git repository with the commend \texttt{git add figures/*}.

\section{Symbols}
Several Latex macros with standard symbols to be used in \belletwo documents are provided in \texttt{belle2-symbols.tex}.

\section{Author list}
Replace the file \texttt{authors.tex} with the actual author list.

\section{References}
The recommended references are
\begin{itemize}
\item \belletwo physics: \cite{Kou:2018nap}
\item B factories: \cite{Bevan:2014iga}
\item Belle: \cite{Brodzicka:2012jm}
\item CKM matrix: \cite{Kobayashi:1973fv}
\item Belle II detector: \cite{Abe:2010gxa}
\item SuperKEKB: \cite{Akai:2018mbz}
\item \evtgen: \cite{Lange:2001uf}
\item \pythia: \cite{Sjostrand:2014zea}
\item \kkmc: \cite{Jadach:1999vf}
\item \tauola: \cite{Jadach:1990mz}
\item \photos: \cite{Barberio:1990ms}
\item \geant: \cite{Agostinelli:2002hh}
\item basf2: \cite{Kuhr:2018lps}
\item FEI: \cite{Keck:2018lcd}
\item PDG: \cite{Zyla:2020zbs}
\item HFLAV: \cite{Amhis:2019ckw}
\end{itemize}
Some other standard references are available in \texttt{references.bib}. To add your own references to the bibtex file use the command \texttt{./addref} with the arXiv number or inspire ID as argument.  The added entries are might still need some editing:
\begin{itemize}
\item check that Latex symbols in the titles are properly rendered (use \belletwo symbols whenever possible);
\item page ranges should in many cases be shortened to just  having the start page;
\item collaborations might need to be added or indeed just the word  ``collaboration'' put after them;
\item for errata and/or addenda use the fields \texttt{extraPrefix}, \texttt{extraVolume}, \texttt{extraPages}, \texttt{extraYear}, and \texttt{extraDoi}, as done for example in Ref.~\cite{Kou:2018nap}.

\end{itemize}

\section{Acknowledgements}
Use the command \texttt{./update} to get the latest version of the acknowledgements and any other updates of the template.

\section{Additional Material}
Add all figures and numbers that should be approved for public presentation in the file \texttt{material.tex}.

\section{Compiling the latex}
Use the \texttt{make} command to compile the latex.
Possible build targets are
\begin{itemize}
\item \texttt{note} (default) for internal notes and for conference notes;
\item \texttt{draft} for paper drafts to be used for collaboration wide review;
\item \texttt{prl} for submission to PRL;
\item \texttt{prd} for submission to PRD;
\item \texttt{jhep} for submission to JHEP;
\item \texttt{epjc} for submission to EPJC.
\end{itemize}

\section{Word counting for PRL}
Use the \texttt{make prlcount} command to perform a word counting following PRL requirements.

\section{Drafting Belle II notes}
Here are examples of questions to ask yourself while drafting a Belle II note.

\subsection*{Basic questions}

\begin{itemize}
\item What are you looking for?
Describe it in one or two sentences, and give a reference.
If you must, you can include a Feynman diagram.

\item What is interesting about it?
Explain the physics motivation for your analysis

\item What previous measurements exist?
Give references to previous measurements and discuss
how your measurement relates to them (\eg more data,
better method, independent cross-check, \etc).

\item How did you look for it?
Keep it very simple: \eg ``we looked for an invariant mass peak in the
 continuum'', or ``we looked for $B$-decays to this final state'', or
``we looked for a two-body $XY$ mass peak in $B \to XYZ$ decays''

\item What data sample did you use?
Give experiment numbers, total int-luminosity;
specify if you used on-resonance data, off-resonance data, or both;
and which skim.

\item Did you use a signal MC sample?
Yes or No;
if Yes, describe the assumptions / generator settings
used to generate it, and the sample size

\item What about background MC?
Yes or No; details as for signal MC

\item What basf2 versions did you use?
Give the basf2 version used for MC generation, skimming, and
analysis (maybe they're different, so be specific), and any other
version-number information that might be relevant here.

\item Where can one find your analysis code?
Give the location of the git repository with your analysis code.
\end{itemize}

\subsection*{Skimming}

\begin{itemize}

\item Did you use a skim?
Apart from using an official skim,
did you use a skimming module to reduce the size of the data/MC?
Answer Yes or No.
If Yes, and it's taken from someone else's analysis, give a reference.
If No, skip the rest of the subsection.

\item If yes, list the cuts used by the skim
An itemized list is fine: \eg invariant mass cuts on combinations,
multiplicity cuts, PID cuts.
If some variable is non-standard, and you also use it in your main
analysis code, you can postpone the description until
the event selection section.

\item What fraction of events pass the skim condition?
If you know the fraction for signal and background MC samples, 
quote that as well.

\end{itemize}

\subsection*{Event selection}

\begin{itemize}

\item How do you select tracks, ECL clusters, \KS candidates \etc for further analysis?
\eg list minimum track momenta and allowed ranges in polar angle; 
 minimum $E_\gamma$ or energy asymmetry cuts for $\pi^0$; 
 V0 cuts for \KS and \Lz; $\pi^0/\eta$ veto for hard gammas...

\item Do you have any cuts against a very general class of background?
 \eg if you're looking for a $B$-decay,
 list your continuum-suppression cuts:
 if they are copied from another Belle II analysis, give the reference.

\item Once you've selected tracks, $\pi^0$ candidates \etc, how do you select events?
 Give an itemized list:
\eg multiplicity cuts, PID cuts,
 valid combinations of particles
 to form your signal candidates, ...

\item How many events survive?
 If you know this for each stage of the analysis (\ie how many events
 survive after each cut) perhaps you could list this in a table,
 especially if you know the answer for signal/background MC as well.
 If not, just write what you do know.

\item You probably have a ``signal region'': how is it defined?
 \eg give $(M_{\rm bc}, \Delta E)$ cuts if appropriate,
 invariant mass windows for your signal candidates...

\item Do you get multiple signal candidates in an event?	If so, how do you handle them?

\item What about sideband regions, or control samples?

\item How many events make it into your signal/sideband regions?
 Again, this should be an easy question. If you know this for signal
 and background MC, also give those numbers here.

\end{itemize}

\subsection*{Data/MC agreement}

\begin{itemize}

\item How did you verify that the MC describes your data?

\item If the MC does not describe the data sufficiently well, how did you determine and apply corrections or calibrations?

\end{itemize}

\subsection*{Analysis}

\begin{itemize}

\item How do you obtain the signal yield?
 \eg using a fit to $\Delta E$ or an invariant mass distribution, 
 or by counting, \etc; give details.

\item How do you make sure your analysis is unbiased?
 Describe the data hiding technique you used.

\item What result do you get (on MC)?
 If you use a fit, show it here, as well as giving the result.
 In any case, show the appropriate plots.
 Keep the signal in data hidden until you got the approval
 to look at it from the reviewers.

\item Are there any other relevant results/plots/\etc?
 \eg show sideband distributions and yields if you know them,
 or yields in control samples...

\item How did you verify that the signal extractions gives a correct estimate of the signal yield?
 Describe how you checked the fit, \eg with pseudo-experiments.

\item How have you estimated the efficiency of your selection/analysis?
 If you're just looking for some final state, perhaps the answer is
 ``I used the MC described at... to estimate it''.
 Or if you're looking for an XY peak in $B \to XYZ$ decays,
 perhaps you've done a check in $M(XY)$ bins.
 Anyway, please give the details.

\item What is the relevant ``physics number'' for this analysis?
 It might be a cross-section; or a product of branching fractions,
 or $\sigma \times\BF\times\BF$;
 or a quasi-two-body fraction
 $\BF(B \to \L Z) \times \BF(\L \to XY) / \BF(B \to XYZ)$.

\item So, what value do you find for this number?
 Please also include all numbers you use in any calculation.

\item Have you done any studies of systematic error? If yes, what studies? With what result?
 Answer Yes or No, describe all details.

\end{itemize}

\subsection*{Other issues}
\begin{itemize}

\item Is there anything important that has been missed so far?

\item Are there other, related analyses being done by you or by people in your group?

\item Do you plan to extend this analysis in the future?
 Basically, answer Yes, No, or Not Sure; and give details if appropriate.

\item Do you have a plan for a longer writeup of this analysis?
 Answer Yes or No; ``No'' is an acceptable answer, if it's the truth.
 If this analysis is part of a wider study (\eg of a particular
 final state in $B$-decay), mention it here.

\end{itemize}